% Standalone: PGR Step-by-Step Chain vs Flat Retrieval (bridge query)
% Compile: pdflatex fig_context.tex
%
% Context format (left column): accurately reconstructed from _bridge_context() in grasp_rag.py
% Answers (both columns): real outputs from data/answers/agriculture_typed.jsonl
\documentclass{article}
\usepackage[paperwidth=22cm, paperheight=16cm,
            top=0.5cm, bottom=0.5cm, left=0.6cm, right=0.6cm]{geometry}
\usepackage{tcolorbox}
\tcbuselibrary{skins,breakable}
\usepackage{xcolor}
\usepackage[T1]{fontenc}
\usepackage{microtype}
\usepackage{amssymb}
\pagestyle{empty}
\setlength{\parindent}{0pt}

% ── Colours ──────────────────────────────────────────────────────────────────
\definecolor{graspblue}{RGB}{37,99,235}
\definecolor{stepbg}{RGB}{239,246,255}
\definecolor{stepborder}{RGB}{147,197,253}
\definecolor{connbg}{RGB}{255,247,237}
\definecolor{connborder}{RGB}{253,186,116}
\definecolor{flatgray}{RGB}{75,85,99}
\definecolor{flatbg}{RGB}{249,250,251}
\definecolor{flatborder}{RGB}{209,213,219}
\definecolor{wincallout}{RGB}{22,163,74}
\definecolor{losecallout}{RGB}{220,38,38}

% ── Box styles ───────────────────────────────────────────────────────────────
\newtcolorbox{entitybox}[1]{
  colback=stepbg, colframe=stepborder, boxrule=0.7pt, arc=3pt,
  top=2pt, bottom=2pt, left=4pt, right=4pt,
  fonttitle=\scriptsize\bfseries\sffamily\color{graspblue},
  title={\strut #1}, titlerule=0pt,
  attach boxed title to top left={yshift=-1.2mm, xshift=3mm},
  boxed title style={colback=white, colframe=stepborder, boxrule=0.5pt,
                     arc=2pt, left=2pt, right=2pt, top=0pt, bottom=0pt},
}
\newtcolorbox{connbox}[1]{
  colback=connbg, colframe=connborder, boxrule=0.7pt, arc=3pt,
  top=2pt, bottom=2pt, left=4pt, right=4pt,
  fonttitle=\scriptsize\bfseries\sffamily\color{orange!70!black},
  title={\strut #1}, titlerule=0pt,
  attach boxed title to top left={yshift=-1.2mm, xshift=3mm},
  boxed title style={colback=white, colframe=connborder, boxrule=0.5pt,
                     arc=2pt, left=2pt, right=2pt, top=0pt, bottom=0pt},
}
\newtcolorbox{flatbox}[1]{
  colback=flatbg, colframe=flatborder, boxrule=0.7pt, arc=3pt,
  top=2pt, bottom=2pt, left=4pt, right=4pt,
  fonttitle=\scriptsize\bfseries\sffamily\color{flatgray},
  title={\strut #1}, titlerule=0pt,
  attach boxed title to top left={yshift=-1.2mm, xshift=3mm},
  boxed title style={colback=white, colframe=flatborder, boxrule=0.5pt,
                     arc=2pt, left=2pt, right=2pt, top=0pt, bottom=0pt},
}

\begin{document}

% ── Header ────────────────────────────────────────────────────────────────────
\begin{center}
  {\small\bfseries\sffamily
   Path-Guided Reasoning (PGR) vs.\ Flat Retrieval — Bridge Query}\\[2pt]
  {\scriptsize\sffamily\color{gray}
   Query: \textit{``How does Varroa Destructor connect to Polyface?''}
   \quad Bridge entity: \textbf{Honey Bees}}
\end{center}

\vspace{3pt}

% ── Two columns ──────────────────────────────────────────────────────────────
\begin{minipage}[t]{0.485\textwidth}

\begin{tcolorbox}[
  colback=graspblue!6, colframe=graspblue, boxrule=1.1pt, arc=4pt,
  top=3pt, bottom=5pt, left=4pt, right=4pt,
  title={\scriptsize\bfseries\sffamily\color{white}
         GRASP-RAG \textnormal{(bridge motif + PGR context)}},
  coltitle=white, attach title to upper,
]

{\tiny\sffamily\color{graspblue!60}\textit{Context sent to LLM
  (built by \texttt{\_bridge\_context()}):}}

\vspace{2pt}

\begin{entitybox}{Step 1 — Entity: Varroa Destructor \normalfont\color{graspblue!60}(parasite)}
\scriptsize\sffamily
Parasitic mite that infests honey bee colonies; causes
parasitic mite syndrome and high colony mortality\ldots\\[-1pt]
\textit{\color{gray!60}[source passage]}
\end{entitybox}

\vspace{2pt}

\begin{connbox}{Step 2 — Connection: Varroa Destructor $\rightarrow$ \textbf{Honey Bees}}
\scriptsize\sffamily
\textbf{Keywords:} colony infestation, bee mortality, mite resistance\\[-1pt]
\textit{\color{gray!60}[source passage]}
\end{connbox}

\vspace{2pt}

\begin{entitybox}{Step 3 — Entity: \textbf{Honey Bees} \normalfont\color{graspblue!60}(livestock)}
\scriptsize\sffamily
Central to pollination of crops; managed at Polyface Farm
as part of integrated sustainable operations\ldots\\[-1pt]
\textit{\color{gray!60}[source passage]}
\end{entitybox}

\vspace{2pt}

\begin{connbox}{Step 4 — Connection: \textbf{Honey Bees} $\rightarrow$ Polyface}
\scriptsize\sffamily
\textbf{Keywords:} pollination services, sustainable agriculture, local food systems\\[-1pt]
\textit{\color{gray!60}[source passage]}
\end{connbox}

\vspace{2pt}

\begin{entitybox}{Step 5 — Entity: Polyface \normalfont\color{graspblue!60}(organisation)}
\scriptsize\sffamily
Joel Salatin's regenerative farm; integrates multiple species
and relies on pollination for crop production\ldots
\end{entitybox}

\vspace{4pt}
{\tiny\sffamily\color{graspblue!60}\textit{Answer generated from context above (real output):}}
\begin{tcolorbox}[colback=wincallout!8, colframe=wincallout!50,
  boxrule=0.7pt, arc=3pt, top=2pt, bottom=2pt, left=4pt, right=4pt]
\scriptsize\sffamily\color{wincallout!80!black}
\textbf{Answer:} ``Varroa Destructor\ldots infests \textbf{honey bee colonies}\ldots
healthy bee populations are essential\ldots for the pollination of
various crops, which Polyface incorporates into its local food systems.''
\quad {\footnotesize\bfseries\checkmark~bridge entity named}
\end{tcolorbox}

\end{tcolorbox}
\end{minipage}
\hfill
\begin{minipage}[t]{0.485\textwidth}

\begin{tcolorbox}[
  colback=flatgray!5, colframe=flatgray!45, boxrule=1.1pt, arc=4pt,
  top=3pt, bottom=5pt, left=4pt, right=4pt,
  title={\scriptsize\bfseries\sffamily\color{white}
         LightRAG / PathRAG \textnormal{(flat retrieval)}},
  coltitle=white, colbacktitle=flatgray!65, attach title to upper,
]

\begin{flatbox}{Vector search results}
\scriptsize\sffamily\color{gray!80}
Passage 1: Varroa Destructor is a major mite pest of honey
bee colonies; causes parasitic mite syndrome\ldots\\[2pt]
Passage 2: Polyface Farm is Joel Salatin's sustainable farm
in Virginia; known for multi-species rotational methods\ldots
\end{flatbox}

\vspace{2pt}

\begin{flatbox}{Keyword graph results}
\scriptsize\sffamily\color{gray!80}
\textit{keywords: sustainable agriculture, pest management,
beekeeping, regenerative farming\ldots}
\end{flatbox}

\vspace{2pt}

{\scriptsize\sffamily\color{gray!60}
\textit{No path between anchors. Varroa Destructor and Polyface
retrieved independently — the connecting node (Honey Bees)
is not labeled or surfaced.}}

\vspace{4pt}

{\tiny\sffamily\color{flatgray!80}\textit{Real outputs from both systems:}}
\begin{tcolorbox}[colback=losecallout!7, colframe=losecallout!40,
  boxrule=0.7pt, arc=3pt, top=2pt, bottom=2pt, left=4pt, right=4pt]
\scriptsize\sffamily\color{losecallout!80!black}
\textbf{LightRAG:} ``I do not have enough information to answer that.''\\[3pt]
\textbf{PathRAG:} ``I'm sorry, but I cannot provide information
connecting Varroa Destructor to Polyface\ldots''\\[3pt]
\quad {\footnotesize\bfseries\color{losecallout}$\times$~bridge entity not found}
\end{tcolorbox}

\vspace{4pt}
\begin{tcolorbox}[colback=yellow!8, colframe=yellow!60!orange,
  boxrule=0.7pt, arc=3pt, top=3pt, bottom=3pt, left=4pt, right=4pt]
\scriptsize\sffamily\color{black!70}
\textbf{Why:} Flat retrieval finds passages \emph{about} each anchor
separately; without a graph path, the connecting node (Honey Bees)
is never retrieved. GRASP-RAG's bidirectional shortest-path
guarantees traversal of the bridge node when both anchors
are in the graph. Win rate on bridge queries vs.\ LightRAG: \textbf{93.4\%}.
\end{tcolorbox}

\end{tcolorbox}
\end{minipage}

\end{document}
